Se calcularon las distancias interplanares del grafito, obteniendo
\(d_{1} = \qty{173.5}{\pico\metre}\) y \(d_{2} = \qty{123.38}{\pico\metre}\)
para el montaje con tubo esférico, mientras que para el montaje con tubo
cilíndrico se obtuvo \(d_{1} = \qty{265.6}{\pico\metre}\) y \(d_{2} = \qty{108.37}{\pico\metre}\).
Además, se determinaron las longitudes de onda para distintos valores de
diferencia de potencial, observándose que esta es inversamente proporcional
a la longitud de onda, en concordancia con el postulado de de Broglie.
Asimismo, se verifico que el diámetro de los círculos concéntricos en el patrón
de difracción es inversamente proporcional a la diferencia de potencial.
Esto se debe a que, al aumentar la energía cinética, los electrones se difractan
en un ángulo menor.
Finalmente, se evidenció el comportamiento ondulatorio de los electrones
a partir el patrón de interferencia generado.
